\documentclass[11pt,letterpaper]{article}

% --- Paquetes Esenciales de Matemática y Tipografía ---
\usepackage[utf8]{inputenc}
\usepackage[spanish,es-tabla]{babel}
\usepackage{amsmath,amssymb,amsfonts,amsthm}
\usepackage{geometry}
\geometry{top=2cm,bottom=2cm,left=2cm,right=2cm}
\usepackage{fancyhdr}
\usepackage{enumitem}
\usepackage{tasks}
\usepackage{tcolorbox}
\tcbuselibrary{skins,breakable}
\usepackage{tabularx}
\usepackage{xcolor}

% --- Paleta de Colores Académica Neutra ---
\definecolor{primaryDark}{RGB}{30, 41, 59}
\definecolor{accentBlue}{RGB}{37, 99, 235}
\definecolor{neutralGray}{RGB}{241, 245, 249}

% --- Configuración de Encabezados y Pies ---
\pagestyle{fancy}
\fancyhf{}
\renewcommand{\headrulewidth}{0.6pt}
\fancyhead[L]{\small\textbf{\textsf{DEPARTAMENTO DE MATEMÁTICA}}}
\fancyhead[R]{\small\textsf{\today}}
\fancyfoot[C]{\small\thepage}

% --- Cajas Didácticas ---
\newtcolorbox{cajainstruccion}{
  colback=neutralGray,
  colframe=primaryDark,
  arc=3mm,
  boxrule=0.8pt,
  fonttitle=\bfseries\sffamily,
  coltitle=white,
  colbacktitle=primaryDark,
  title=Instrucciones Generales
}

\newtcolorbox{cajaejemplo}[1]{
  colback=white,
  colframe=accentBlue,
  arc=2mm,
  boxrule=1pt,
  fonttitle=\bfseries\sffamily,
  title=#1
}

\begin{document}

% --- Membrete Estándar Institucional ---
\noindent
\begin{tabularx}{\textwidth}{|X|l|l|}
\hline
\textbf{Estudiante:} & \textbf{Curso:} & \textbf{Fecha:} \\
\hline
\textbf{Evaluación:} Evaluación Formativa: Función y Ecuación Cuadrática en Contexto & \textbf{Puntaje Total:} 12 pts & \textbf{Nota:} \\
\hline
\end{tabularx}

\vspace{0.4cm}

\begin{cajainstruccion}
Lea atentamente cada situación problema. Justifique formalmente sus procedimientos matemáticos. En los ítems de selección múltiple, marque únicamente la opción correcta.
\end{cajainstruccion}

\vspace{0.4cm}

% --- Cuerpo Generado Dinámicamente por el Formatter ---
\noindent \textbf{Ítem 1} (3 pts)\\
Considere la ecuación cuadrática $x^2 - 6x + 8 = 0$. ¿Cuál es el conjunto solución en $\mathbb{R}$?

\begin{tasks}(2)
  \task \textbf{A)} x = 2 y x = 4
  \task \textbf{B)} x = -2 y x = -4
  \task \textbf{C)} x = 1 y x = 8
  \task \textbf{D)} x = 3 y x = -3
\end{tasks}

\vspace{0.4cm}

\noindent \textbf{Ítem 2} (4 pts)\\
La trayectoria de un proyectil lanzado verticalmente está modelada por la función $h(t) = -5t^2 + 20t$, donde $h$ es la altura en metros y $t$ el tiempo en segundos ($t \geq 0$). ¿En qué instante $t$ el proyectil alcanza su altura máxima?

\begin{tasks}(2)
  \task \textbf{A)} t = 2 segundos
  \task \textbf{B)} t = 4 segundos
  \task \textbf{C)} t = 20 segundos
  \task \textbf{D)} t = 5 segundos
\end{tasks}

\vspace{0.4cm}

\noindent \textbf{Ítem 3} (5 pts)\\
Un estudiante afirma que 'si el discriminante $\Delta = b^2 - 4ac$ de una ecuación cuadrática es estrictamente negativo, la ecuación no tiene solución'. Demuestre o refute dicha afirmación especificando con rigor los conjuntos numéricos (conjunto $\mathbb{R}$ vs. $\mathbb{C}$) y su interpretación geométrica en el plano cartesiano.

\vspace{0.2cm}
\begin{tcolorbox}[colback=white,colframe=neutralGray,arc=1mm,boxrule=0.5pt,height=4.5cm,valign=top]
{\footnotesize\color{gray}\textit{Espacio reservado para desarrollo y justificación axiomática}}
\end{tcolorbox}

\vspace{0.4cm}


\end{document}
